Figure~\ref{fig:result_widths} shows the final measured widths of 
  the away-side correlation (standard-deviation of $Y^{\mathrm{corr}}(\Dphi)$ in Eq.~\eqref{eq:fit_func})
  in \PbPb\ events, as a function of centrality.
The value measured in \pp\ collisions is also shown.
The measurements are shown for All pairs, Same-sign pairs 
  and Opposite-sign pairs.
The widths of the same- and opposite-sign pairs are similar 
  with each other for both the \PbPb\ and \pp\ measurements.
The measured \PbPb\ values are very similar to the values measured in 
  \pp\ collisions.
No clear increase in the widths as a function of centrality 
  is observed.
Instead a slight decrease in the 0--10\% most central events is observed.
This is very surprizing as one would expect that interactions
  with the QGP would broaden the \Dphi\ correlations~\cite{Nahrgang:2013saa}.
This result thus indicates that while the HF quarks are strongly quenched,
  they do not undergo significant angular deflection in the QGP. 
Figure~\ref{fig:result_widths_fits} (\ref{fig:result_widths_const_fits})
  shows linear (constant) fits to the
  \PbPb\ widths as a function of centrality.
The slopes are negative for all three charge-combinaitons
  but consistent with zero within 3$\sigma$ statistical uncertainties.
Similar fits but with the 0--10\% centrality interval excluded 
  are shown in Appendix~\ref{sec:fits_slope}.
Those plots show that over the 10--80\% centrality range,
  there is very little change in the widths.

Using these measurements, we can put constraints on the \Dphi\ 
  broadening introduced by the QGP interactions.
Figure~\ref{fig:result_widths_diff} shows the squared difference
  between the \pp\ and \PbPb\ widths.
Figure~\ref{fig:result_widths_ratio} shows the ratios of the widths.
For these plots, correlations in the systematic uncertainties
  in the \PbPb\ and \pp\ measurements are accounted for as follows:
  for each Systematic variation (discussed in 
  Section~\ref{sec:systematics} and Appendix~\ref{sec:systematicsPP}),
  the \PbPb/\pp\ ratio (or r.m.s. difference) for the widths are evaluated. 
Then the spread in the ratios from each source from the nominal ratio
  are added in quadrature to get the uncertainty of the ratio.
Except for in the 0--10\% centrality interval,
  the \PbPb\ values, are within $\sim$1~$\sigma$
  statistical/systematic uncertainties to the \pp\ values.
Figure~\ref{fig:result_widths_diff} also shows the 90\% C.L.
  on the upper limit of the difference, evaluated using the
  combined statistical + systematin uncertainties.
The lower panels in Figure~\ref{fig:result_widths_diff}
  also show the estimated upper-limit on the 
  broadening introduced by the interactions with the QGP:
\begin{eqnarray}
\sigma_{\mathrm{int}}=\sqrt{\sigma^{2}_{\mathrm{Pb+Pb}} - \sigma^{2}_{pp}}.
\label{eq:broaden}
\end{eqnarray}
The square-root in Eq.~\eqref{eq:broaden}, is negative for the 90\% C.L. value 
  for the 0--10\% centrality interval and is not shown 
  in Figure~\ref{fig:result_widths_diff}.
  

\begin{figure}
\centering
\includegraphics[width=1.0\textwidth]%                                                                                           
{figures/AllFigs/can_systPbPbPP_WidthTau_Tight_pTBar_GT5_EffCor}
\caption{
The measured widths of the away-side correlation.
The plot shows the $\sigma$ (standard-deviation of $Y^{\mathrm{corr}}(\Dphi)$ in Eq.\eqref{eq:fit_func})
  as well as the $\Gamma$ (the half-width at half-maximum parameter in Eq.\eqref{eq:fit_func}))
  for dimuon pairs, as a function of centrality.
Also shown for comparison, are the corresponding parameters in \pp\ collisions (solid markers).
  The horizontal lines indicates the nominal \pp\
  values (i.e. the central value of the \pp\ measurements,
  but without the statistical or systematic uncertainties) drawn across the full centrality range.
The left, middle and right panels correspond to
  Same-sign pairs, Opposite-sign pairs and All-pairs.
The vertical lines and shaded bars indicate statistical and
  systematic uncertainties, respectively.
\label{fig:result_widths}
}
\end{figure}



\begin{figure}
\centering
\includegraphics[width=1.0\textwidth]%
{figures/AllFigs/can_fits_0_Tight_pTBar_GT5_EffCor}
\caption{
\label{fig:result_widths_fits}
Linear fits to the measured widths of the away-side correlation.
The left, middle and right panels correspond to
  Same-sign pairs, Opposite-sign pairs and All-pairs.
The top row shows $\chi^2$-fits  made taking into consideration the statistical uncertainties only.
The middle (bottom) row shows $\chi^2$-fits  made taking into consideration the systematic  uncertainties only
  while treating the uncertainties fully uncorrelated (correlated) bin-to-bin.
The fit values and their uncertainties are shown on the individual panels.
}
\end{figure}

\begin{figure}
\centering
\includegraphics[width=1.0\textwidth]%
{figures/AllFigs/can_constant_fits_0_Tight_pTBar_GT5_EffCor}
\caption{
\label{fig:result_widths_const_fits}
Same as Figure~\ref{fig:result_widths_fits} but using fits to Constant-functions.
}
\end{figure}

\begin{figure}
\centering
\includegraphics[width=1.0\textwidth]%
{figures/AllFigs/can_Rmsdiff_New_0_Tight_pTBar_GT5_EffCor}
\includegraphics[width=1.0\textwidth]%
{figures/AllFigs/can_SmearingSigmaCL_New_0_Tight_pTBar_GT5_EffCor}
\caption{
Top panel: The difference between the squared-widths of the \pp\ and \PbPb\ measurements.
The left, middle and right panels correspond to
  Same-sign pairs, Opposite-sign pairs and All-pairs.
The vertical lines and shaded bars indicate statistical and
  systematic uncertainties, respectively.
The uncertainty accounts for the correlated systematic uncertainties
  between the \PbPb\ and \pp\ measurements.
The blue lines show the 90\% C.L. upper limit on the difference.
Bottom panels: The 90\% C.L. on the broadening 
  introduced by the interactions of the muons in the QGP (see text).
\label{fig:result_widths_diff}
}
\end{figure}


\begin{figure}
\centering
\includegraphics[width=1.0\textwidth]%
{figures/AllFigs/can_ratio_0_Tight_pTBar_GT5_EffCor}
\caption{
The ratio between the squared-widths of the \PbPb\ and \pp\ measurements.
The left, middle and right panels correspond to
  Same-sign pairs, Opposite-sign pairs and All-pairs.
The vertical lines and shaded bars indicate statistical and
  systematic uncertainties, respectively.
The uncertainty accounts for the correlated systematic uncertainties
  between the \PbPb\ and \pp\ measurements.
\label{fig:result_widths_ratio}
}
\end{figure}


\iftoggle{IncludeYields}
{
  Figure~\ref{fig:result_yields}, shows the measured yields of 
    the away-side correlation (integral of $Y^{\mathrm{corr}}(\Dphi)$ in Eq.\eqref{eq:fit_func})
    in \PbPb\ events, as a function of centrality.
  Figure~\ref{fig:result_yields_taa}, shows similar plots, but with for the
    per-event yields scaled by the \TAA (Eq.~\eqref{eq:TaaYield}).
  The cross-section measured in \pp\ collisions is also shown in Figure~\ref{fig:result_yields_taa}.
  As expected, the yields (Figure~\ref{fig:result_yields}) increase monotonically
    from peripheral to central events.
  However the \TAA\ scaled yields show a monotonic decrease. 
  This indicates that there is a significant suppression in the number of dimuon pairs
    compared to what one would expect if the \PbPb\ collisions were
    a superimposition of independent \pp\ collisions; 
    in which case the scaled yields would be a constant and consistent 
    with the cross-section measured in \pp\ collisions.
  The relative decrease in the \TAA\ scaled-yields are similar for the same-sign pairs 
    and the opposite-sign pairs.

  \begin{figure}
  \centering
  \includegraphics[width=1.0\textwidth]%                                                                                           
  {figures/AllFigs/can_syst_1_Tight_pTBar_GT5_EffCor}
  \caption{
  The measured correlated yields on the away-side 
    (integral of $Y^{\mathrm{corr}}(\Dphi)$ in Eq.\eqref{eq:fit_func}) for dimuon pairs, 
    as a function of centrality.
  The left, middle and right panels correspond to
    All-pairs, Same-sign pairs and Opposite-sign pairs.
  The vertical lines and shaded bars indicate statistical and 
    systematic uncertainties, respectively.
  The overall luminosity uncertainty of $\pm1.5$\% is not included
    in the shaded bands, instead stated on the plot.
  \label{fig:result_yields}
  }
  \end{figure}

  \begin{figure}
  \centering
  \includegraphics[width=1.0\textwidth]%                                                                                           
  {figures/AllFigs/can_systPbPbPP_2_Tight_pTBar_GT5_EffCor}
  \caption{
  The \TAA\ scaled yields on the away-side for dimuon pairs, 
    as a function of centrality.
  Also shows for comparison, is the \pp\ cross-section (blue marker).
  The horizontal blue-line indicates the nominal \pp\
    value drawn across the full centrality range.
  The left, middle and right panels correspond to
    All-pairs, Same-sign pairs and Opposite-sign pairs.
  The vertical lines and shaded bars indicate statistical and 
    systematic uncertainties, respectively.
  The overall luminosity uncertainties of $\pm1.5$\% for \PbPb\
    and $\pm1.6$\% for \pp\ are not included
    in the shaded bands, instead stated on the plot.
  \label{fig:result_yields_taa}
  }
  \end{figure}
}
{}

\clearpage
\section{Summary}
This note reports a measurement of correlated back-to-back 
  muon-pair production in \PbPb\ collisions at \snn=5.02~\TeV. 
  and in \pp\ collisions at \sqs=5.02~\TeV.
The \PbPb\ measurements are performed using data from the 2015 and 2018
  heavy-ion runs corresponding to an integrated luminosity of
  1.93~\inb.
While the \pp\ measurements are performed using data from the 2017
  reference run corresponding to an integrated luminosity of 0.26~\ifb.
The measurements are performed over the integrated \ptbar>5~\GeV\ range.
Over this \pt\ range, nearly all the dimuon pairs arise from decay
  of HF hadrons.
Backgrounds from the underlying event (including flow-correlated pairs) 
  are removed.
Backgrounds from muons arising from inflight decays of pions and Kaons,
  as well as muons arising from secondary interactions in the calorimeter 
  were estimated using the momentum-imbalance as a discriminant, 
  and were found to be below a few percent.
\iftoggle{IncludeYields}
{
  The yields of the correlated muon-pairs, as well as the width of the
    correlation in \Dphi\ are measured.
  The per-event yields, scaled by the nuclear thickness function, \TAA, 
    are observed to decrease from peripheral to central collisions,
    indicating a suppression in the yield of such pairs in central collisions.
  On the other hand, the widths of the correlations are observed to be 
    quite independent of centrality.
  Theoretical models~\cite{Nahrgang:2013saa} predict varying amounts of 
    suppression and broadening of back-to-back HF pairs depending 
    on the energy-loss mechanism.
  Thus these measurements can shed light on the energy-loss mechanism of 
    HF-quarks in the QGP produced in heavy-ion collisions.
}
{
  The width of the correlated muon-pairs in \Dphi\ are measured,
    and are observed to be independent of centrality, within uncertainties.
  The widths are also observed to be consistent between the \pp\ and the \PbPb\ systems.
  Theoretical models~\cite{Nahrgang:2013saa} predict varying amounts of 
    centrality-dependent broadening of such back-to-back HF pairs in \PbPb\ collisions
    depending on the heavy-quark energy-loss mechanism.
  Thus these measurements can shed light on the energy-loss mechanism of 
    HF-quarks in the QGP produced in heavy-ion collisions.
}
